\scanpage{217}

\diagram{\begin{tikzcd}[ampersand replacement=\&,column sep=small,row sep=small]
X\arrow[r,"{\unsure{F|f}}_{X/S}"]\arrow[d,dashed,dash]\&X^{p/S}\arrow[r,dash]\arrow[dl]\&X\arrow[dl]\\
S\arrow[r,dash]\&S
\end{tikzcd}}

Considérons $F_{X/S}^*(C\omega)\in\Gamma F_{X/S}^*\underline\Omega^1_{X^{p/S}/S}\simeq\Gamma\underline\Omega^1_{X/S}{}^{(p)}$, on a alors

\[
\begin{aligned}
u'(\xi)&=u(\xi)+\langle\omega^{\otimes p}-F_{X/S}^*(C\omega),\xi^{\otimes p}\rangle,\quad\text{i.e. intrinsèquement}\\
u'&=u+\underbrace{(\omega^{\otimes p}-F_{X/S}^*C\omega)}_{F^*_{X/S}(\widetilde\omega-C\omega)}
\end{aligned}
\]

Donc la nouvelle obstruction se trouve dans

\[
\Gamma(X,\Omega^1_{X/S}{}^{\otimes p})/\mathrm{Im}\,[Z^1(X,\Omega^*_{X/S})\to\Gamma(X,\Omega^1_{X/S}{}^{\otimes p})]
\]

l'hom. étant
\[
Z^1(X,\Omega^*_{X/S})\longrightarrow\Gamma(X,\Omega^1_{X/S}{}^{\otimes p})\qquad
\omega\longmapsto F_{X/S}^*(\widetilde\omega-C\omega)=\omega^{\otimes p}-F_{X/S}^*(C\omega)
\]

\margin{$2$}

Supposons $X$ propre sur $S$, on trouve, en localisant sur $S$, un foncteur

\[
\deleted{f_*(\underline Z^1_{X/S})}\,
\underset{\deleted{\substack{\|\\\text{faisceau des}\\\text{inversibles à}\\\text{connexions intégrables}\\\text{de }X/S}}}{\deleted{\underline{\mathrm{Pic}}_{X/S}^{\unsure{c}}}}
\inserted{\underline{\mathrm{Pic}}\supset\underline{\mathrm{Pic}}^{\text{conn int}}}
\longrightarrow \deleted{f_*\Omega^{1\otimes p}}\,
[\,\underset{\deleted{\unsure{C|R}}}{\mathrm{Coker}}\,\underbrace{f_*\,\underline Z^1_{X/S}}_{\|\,\mathrm{Ker}[f_*(\Omega^1_{X/S})\xrightarrow d f_*(\Omega^2_{X/S})]}
\xrightarrow{\gamma}f_*(\Omega^{1\otimes p}_{X/S})\,]\tag*{(1)}
\]
\margin{pas linéaire !}

Supposons que

a) $\underline\Omega^1_{X/S},\underline\Omega^2_{X/S},\Omega^1_{X/S}{}^{\otimes p}$ sont coh. plats sur $S$ en dim 0

b) $\underline{\mathrm{Co}}\mathrm{ker}\ d^{1}$ et $\mathrm{Im}\,d^1$ sont plats sur $S$, où $d^1:f_*(\underline\Omega^1_{X/S})\to f_*(\underline\Omega^2_{X/S})$ [OK si $X$ de dim relative 1 sur $S$]

c) \deleted{\unsure{C|le}} $\mathrm{Ker}\,\gamma$ est plat sur $S$.

Alors le 2\textsuperscript{ème} mbr de \unsure{(1)|(3)|(ci)} est représentable par un fibré vectoriel

% archive sheet number: 216
